\documentclass[a4paper,11pt]{article}

\usepackage{revy}
\usepackage[utf8]{inputenc}
\usepackage[T1]{fontenc}
\usepackage[danish]{babel}
\usepackage{amsmath,amssymb,amsfonts,amsthm,textcomp}


\revyname{MatematikRevyen}
\revyyear{2026}
% HUSK AT OPDATERE VERSIONSNUMMER
\version{0.1}
\eta{$2$ minutter}
\status{Færdig}

\title{Løvens A-lokale 2}
\author{Carl '21 \& Thais '22}

\begin{document}
\maketitle

\begin{roles}
\role{E}[Skuespiller] Studerende-entrepeneur (forvirret type)
\role{S}[Skuespiller] Speaker (meget få linjer, ku være i bandet?)
\role{L1}[Skuespiller] Forelæser-løve (Ernst)
\role{L2}[Skuespiller] Forelæser-løve (Morten)
\role{L3}[Skuespiller] Forelæser-løve 
\end{roles}

\begin{props}
\item Bord med tre stole til løverne
\end{props}


\begin{sketch}
\says{S} Den næste iværksætter er E som ikke dukket op endnu

\says{L1} Det fandme uprof. det der. 
\says{L2} Ja, hvem kommer ikke til tiden?

\sasy{E} Hej løver, i er godt nok mødt tidligt ind. Hvad ville I sige til at kunne komme for sent uden rent faktisk at komme for sent?

\scene{Intro-jingle}

\says{E} Forestil dig, at der står du skal møde kl  10:00, hvad tid dukker du op? 

\says{L2} Jeg ville nok være der 5 min før

\says{E} Men hvad hvis du kunne møde 10min efter og stadig være i god tid?!

\says{L2 Det kan man da ikke?

\says{E} Det kan man godt, og det er kernen i mit koncept - Akademisk kvarter. Med akademisk kvarter indfører vi en helt ny måde at tænke tid på, hvor du kan møde kvart over 10 selvom din forelæsning starter kl 10.

\says{L3} Men starter forelæsningen så kl 10?

\says{E} Nej eller jo, både og. Men du dukker i hvert fald op kvart over. Forstår i?

\says{L1} Skal vi så fortælle folk, at vi bruger akademisk kvarter

\says{E} Nej, det ville være indforstået. Jeg vil fremelske en kultur hvor alle møder ind et kvarter efter vi fortæller dem de skal.

\says{L3} Jeg forstår simpelthen ikke hvor vi ikke bare beder folk om at møde kl 10:15 i stedet

\says{E} Fordi så kommer de jo for sent 

\says{L1} Men forstår jeg ret, at du gerne vil have at folk kommer for sent?

\says{E} Ja det er simpelthen indbygget i den måde vi lever. Så folk kommer for sent, men ikke for sent for sent, I ved.

\says{L2} Jeg ved overhovedet ikke noget lige nu. Der er en rigtig og en forkert måde at komme for sent på?

\says{L1} Okay jeg tror jeg har den. Så når jeg har undervisning i LinAlg kl 8:00, så starter selve undervisningen vi kl 8:15

\says{E} Nej! Lige der, der starter du kl 8 sharp. 

\scene{Lys ned}


\end{sketch}
\end{document}
